\chapter{Dokumentasi Praktikum}

\section{Hasil Praktikum}

\subsection{Inisiasi Proyek dan Pembuatan Blueprint}

\begin{figure}[H]
    \centering
    \includegraphics[width=0.6\textwidth]{1.png}
    \caption{Pembuatan Proyek Baru}
    \label{fig:1}
\end{figure}
Keterangan Gambar :
\begin{itemize}
    \item Membuat proyek baru di Unreal Engine menggunakan template \textbf{Third Person}.
    \item Penamaan proyek disesuaikan dengan format yang telah ditentukan sebelum menekan tombol \textbf{Create}.
\end{itemize}

\begin{figure}[H]
    \centering
    \includegraphics[width=0.6\textwidth]{2.png}
    \caption{Menambahkan Blueprint Class}
    \label{fig:2}
\end{figure}
Keterangan Gambar :
\begin{itemize}
    \item Melakukan klik kanan pada \textit{Content Drawer} untuk membuat \textit{Blueprint Class} baru.
\end{itemize}

\begin{figure}[H]
    \centering
    \includegraphics[width=0.6\textwidth]{3.png}
    \caption{Pemilihan Parent Class}
    \label{fig:3}
\end{figure}
Keterangan Gambar :
\begin{itemize}
    \item Memilih \textbf{Actor} sebagai \textit{Parent Class} dari Blueprint yang dibuat.
    \item Mengubah nama Blueprint menjadi \textbf{BP\_BOX} sesuai dengan format yang diinstruksikan.
\end{itemize}

\begin{figure}[H]
    \centering
    \includegraphics[width=0.6\textwidth]{4.png}
    \caption{Tampilan Viewport pada Blueprint}
    \label{fig:4}
\end{figure}
Keterangan Gambar :
\begin{itemize}
    \item Setelah Blueprint \textbf{BP\_BOX} dibuka, antarmuka akan menampilkan bagian \textbf{Viewport} yang digunakan untuk mengatur komponen visual dari \textit{Actor} tersebut.
\end{itemize}

\begin{figure}[H]
    \centering
    \includegraphics[width=0.6\textwidth]{5.png}
    \caption{Tampilan Event Graph pada Blueprint}
    \label{fig:5}
\end{figure}
Keterangan Gambar :
\begin{itemize}
    \item Beralih ke tab \textbf{Event Graph}, yang merupakan tempat untuk menyusun logika visual (\textit{node-based programming}) dari Blueprint.
\end{itemize}

\subsection{Eksplorasi Node dan Fungsi pada Blueprint}

\begin{figure}[H]
    \centering
    \includegraphics[width=0.6\textwidth]{6.png}
    \caption{Menambahkan Variabel Baru}
    \label{fig:6}
\end{figure}
Penjelasan :
\begin{itemize}
    \item \textbf{Cara:} Pada panel \textit{My Blueprint} di sebelah kiri, klik tombol \textbf{+} pada bagian \textit{Variables}.
    \item \textbf{Fungsi:} Digunakan untuk menyimpan data atau nilai tertentu (seperti \textit{Boolean, Integer, Float}, dll.) yang dapat digunakan dan diubah selama permainan berlangsung.
\end{itemize}

\begin{figure}[H]
    \centering
    \includegraphics[width=0.6\textwidth]{7.png}
    \caption{Node Action: Print String}
    \label{fig:7}
\end{figure}
Penjelasan :
\begin{itemize}
    \item \textbf{Cara:} Klik kanan pada \textit{Event Graph} dan cari \textbf{Print String}.
    \item \textbf{Fungsi:} Berfungsi untuk menampilkan pesan teks ke layar saat permainan dijalankan. Sangat berguna untuk proses \textit{debugging} dan mengecek apakah suatu logika berjalan dengan benar.
\end{itemize}

\begin{figure}[H]
    \centering
    \includegraphics[width=0.6\textwidth]{8.png}
    \caption{Menambahkan Structure Baru}
    \label{fig:8}
\end{figure}
Penjelasan :
\begin{itemize}
    \item \textbf{Cara:} Melalui \textit{Content Drawer}, klik kanan dan pilih \textit{Blueprints}, lalu pilih \textbf{Structure}.
    \item \textbf{Fungsi:} Membuat kumpulan berbagai tipe variabel berbeda yang dibungkus menjadi satu kesatuan data. Berguna untuk mengorganisir data kompleks (misal: status pemain berisi nyawa, mana, dan kecepatan).
\end{itemize}

\begin{figure}[H]
    \centering
    \includegraphics[width=0.6\textwidth]{9.png}
    \caption{Tampilan Editor Structure}
    \label{fig:9}
\end{figure}
Penjelasan :
\begin{itemize}
    \item \textbf{Cara:} Buka \textit{Structure} yang telah dibuat sebelumnya dengan melakukan klik dua kali.
    \item \textbf{Fungsi:} Antarmuka ini menampilkan daftar variabel di dalam \textit{Structure}, di mana kita bisa menambahkan variabel baru, mengubah tipe data, dan menentukan nilai bawaannya (\textit{default value}).
\end{itemize}

\begin{figure}[H]
    \centering
    \includegraphics[width=0.6\textwidth]{10.png}
    \caption{Node Action: Branch}
    \label{fig:10}
\end{figure}
Penjelasan :
\begin{itemize}
    \item \textbf{Cara:} Klik kanan pada \textit{Event Graph}, cari \textbf{Branch} (atau tekan dan tahan tombol 'B' lalu klik kiri).
    \item \textbf{Fungsi:} Node percabangan logika berdasarkan kondisi \textit{Boolean} (Benar/Salah). Jika kondisi yang dievaluasi bernilai \textit{True}, alur eksekusi akan mengalir ke pin \textit{True}, begitu pula sebaliknya untuk \textit{False}.
\end{itemize}

\begin{figure}[H]
    \centering
    \includegraphics[width=0.6\textwidth]{11.png}
    \caption{Memasukkan Variabel ke Event Graph (Getter dan Setter)}
    \label{fig:11}
\end{figure}
Penjelasan :
\begin{itemize}
    \item \textbf{Cara:} Lakukan metode \textit{drag and drop} dari panel variabel ke dalam lembar kerja \textit{Event Graph}. Saat diseret dan ditahan, akan muncul penanda arah kursor variabel.
    \item \textbf{Fungsi:} Ketika dilepas (\textit{drop}), akan muncul pilihan untuk menjadikannya \textbf{Get} (mengambil/membaca nilai variabel) atau \textbf{Set} (mengubah/menulis nilai baru pada variabel tersebut). Kedua simulasi visual ini sangat penting dalam pengaturan data logika.
\end{itemize}

\begin{figure}[H]
    \centering
    \includegraphics[width=0.6\textwidth]{12.png}
    \caption{Node Action: Event BeginPlay}
    \label{fig:12}
\end{figure}
Penjelasan :
\begin{itemize}
    \item \textbf{Cara:} Node ini umumnya sudah ada secara bawaan, namun dapat dipanggil kembali dengan mengetik \textbf{Event BeginPlay}.
    \item \textbf{Fungsi:} Sebagai pemicu (\textit{trigger}) awal. Segala logika yang terhubung dengan node ini akan dieksekusi tepat satu kali ketika permainan (atau \textit{Actor} ini) pertama kali dimuat dan dimulai.
\end{itemize}

\begin{figure}[H]
    \centering
    \includegraphics[width=0.6\textwidth]{13.png}
    \caption{Node Action: Event Tick}
    \label{fig:13}
\end{figure}
Penjelasan :
\begin{itemize}
    \item \textbf{Cara:} Sama seperti BeginPlay, node ini sering tersedia dari awal atau bisa dicari dengan kata kunci \textbf{Event Tick}.
    \item \textbf{Fungsi:} Memicu jalannya eksekusi secara terus-menerus pada setiap \textit{frame} (gambar) permainan. Node ini juga menyediakan output \textit{Delta Seconds} yang berfungsi untuk sinkronisasi waktu pergerakan independen dari kecepatan \textit{frame rate}.
\end{itemize}

\begin{figure}[H]
    \centering
    \includegraphics[width=0.6\textwidth]{14.png}
    \caption{Node Action: Get Player Controller}
    \label{fig:14}
\end{figure}
Penjelasan :
\begin{itemize}
    \item \textbf{Cara:} Cari node dengan mengetik \textbf{Get Player Controller}.
    \item \textbf{Fungsi:} Node ini mengambil referensi \textit{Controller} yang sedang dikendalikan oleh pemain. Berguna sebagai jembatan untuk mengelola pengaturan input pemain atau antarmuka (UI).
\end{itemize}

\begin{figure}[H]
    \centering
    \includegraphics[width=0.6\textwidth]{15.png}
    \caption{Node Action: Enable Input}
    \label{fig:15}
\end{figure}
Penjelasan :
\begin{itemize}
    \item \textbf{Cara:} Cari node dengan kata kunci \textbf{Enable Input}.
    \item \textbf{Fungsi:} Mengizinkan \textit{Actor} tersebut (misal: pintu, kendaraan) untuk dapat menerima dan mendengarkan masukan (input) keyboard atau mouse dari pemain secara langsung, dengan menghubungkan referensi \textit{Player Controller}.
\end{itemize}

\begin{figure}[H]
    \centering
    \includegraphics[width=0.6\textwidth]{16.png}
    \caption{Node Action: Keyboard Event}
    \label{fig:16}
\end{figure}
Penjelasan :
\begin{itemize}
    \item \textbf{Cara:} Cari node dengan mengetik tombol spesifik yang diinginkan (contoh: \textbf{Keyboard R}).
    \item \textbf{Fungsi:} Merupakan pemicu aksi yang bereaksi ketika sebuah tombol ditekan (\textit{Pressed}) atau dilepas (\textit{Released}). Node ini tidak akan bekerja di dalam \textit{Actor} jika fungsi input belum diaktifkan (\textit{Enable Input}).
\end{itemize}

\begin{figure}[H]
    \centering
    \includegraphics[width=0.6\textwidth]{17.png}
    \caption{Node Action: Branch (Lanjutan)}
    \label{fig:17}
\end{figure}
Penjelasan :
\begin{itemize}
    \item \textbf{Deskripsi Tambahan:} Ini merupakan tampilan penggunaan \textit{Branch} dalam kombinasi fungsi lain.
    \item \textbf{Fungsi:} Mengevaluasi sebuah nilai Boolean. Sangat vital untuk menyaring suatu kondisi sebelum mengeksekusi aksi lanjutan, misalnya mengecek apakah pemain memiliki kunci pintu atau tidak.
\end{itemize}

\begin{figure}[H]
    \centering
    \includegraphics[width=0.6\textwidth]{18.png}
    \caption{Node Action: Print String (Lanjutan)}
    \label{fig:18}
\end{figure}
Penjelasan :
\begin{itemize}
    \item \textbf{Deskripsi Tambahan:} Variasi pemasangan \textit{Print String} dengan input \textit{String} kustom.
    \item \textbf{Fungsi:} Teks \textit{String} yang ditampilkan dapat dimodifikasi durasi dan warnanya, sehingga \textit{programmer} bisa dengan mudah membedakan hasil eksekusi dari beberapa pesan \textit{debug} yang muncul bersamaan.
\end{itemize}

\begin{figure}[H]
    \centering
    \includegraphics[width=0.6\textwidth]{19.png}
    \caption{Membuat Variabel Bertipe Boolean}
    \label{fig:19}
\end{figure}
Penjelasan :
\begin{itemize}
    \item \textbf{Cara:} Saat membuat variabel, pastikan untuk mengganti tipe variabel menjadi \textbf{Boolean} (ditandai dengan warna merah).
    \item \textbf{Fungsi:} Tipe data ini hanya menyimpan dua nilai, yaitu \textit{True} atau \textit{False}. Umumnya digunakan sebagai bendera status (\textit{flag}) atau parameter kondisi untuk node \textit{Branch}.
\end{itemize}

\begin{figure}[H]
    \centering
    \includegraphics[width=0.6\textwidth]{20.png}
    \caption{Node Action: Do Once}
    \label{fig:20}
\end{figure}
Penjelasan :
\begin{itemize}
    \item \textbf{Cara:} Ketik dan panggil node \textbf{Do Once}.
    \item \textbf{Fungsi:} Sesuai namanya, node ini membatasi aliran eksekusi agar hanya dapat dilewati sebanyak \textbf{satu kali} saja. Eksekusi lanjutan akan ditolak kecuali ada sinyal baru yang masuk dari jalur pin \textit{Reset}.
\end{itemize}

\begin{figure}[H]
    \centering
    \includegraphics[width=0.6\textwidth]{21.png}
    \caption{Node Action: Flip Flop}
    \label{fig:21}
\end{figure}
Penjelasan :
\begin{itemize}
    \item \textbf{Cara:} Ketik dan panggil node \textbf{Flip Flop}.
    \item \textbf{Fungsi:} Membagi eksekusi yang masuk secara bergantian (bolak-balik). Panggilan pertama dialirkan ke pin \textbf{A}, panggilan kedua dialirkan ke pin \textbf{B}, panggilan ketiga kembali ke \textbf{A}, dan seterusnya. Biasa digunakan untuk tombol pemicu (\textit{toggle}), misalnya menghidupkan/mematikan lampu.
\end{itemize}

\begin{figure}[H]
    \centering
    \includegraphics[width=0.6\textwidth]{22.png}
    \caption{Node Action: Gate}
    \label{fig:22}
\end{figure}
Penjelasan :
\begin{itemize}
    \item \textbf{Cara:} Ketik dan panggil node \textbf{Gate}.
    \item \textbf{Fungsi:} Bertindak sebagai gerbang logika. Sinyal yang masuk dari pin \textit{Enter} hanya bisa diteruskan ke \textit{Exit} asalkan gerbang dalam keadaan \textit{Open}. Jika gerbang tertutup (\textit{Close}), eksekusi akan tertahan. Memiliki pin pengatur seperti buka, tutup, dan saklar otomatis (\textit{Toggle}).
\end{itemize}

\begin{figure}[H]
    \centering
    \includegraphics[width=0.6\textwidth]{23.png}
    \caption{Node Action: Switch}
    \label{fig:23}
\end{figure}
Penjelasan :
\begin{itemize}
    \item \textbf{Cara:} Ketik dan panggil node \textbf{Switch} (sesuai parameter seperti \textit{Switch on Int}, \textit{Switch on Enum}, dll).
    \item \textbf{Fungsi:} Mirip dengan \textit{Branch} namun dapat mengevaluasi lebih dari dua kondisi (memiliki banyak jalur percabangan). Node ini mengalihkan rute eksekusi berdasarkan nilai variabel dari parameter yang dievaluasi.
\end{itemize}
\subsection{Implementasi Logika Blueprint dan Hasil Akhir}

\begin{figure}[H]
    \centering
    \includegraphics[width=0.6\textwidth]{24.png}
    \caption{Implementasi Blueprint: Event BeginPlay}
    \label{fig:24}
\end{figure}
Penjelasan :
\begin{itemize}
    \item \textbf{Cara Kerja:} Skenario percakapan logika paling dasar. Saat permainan dijalankan, pemicu \textit{Event BeginPlay} akan menyala dan mengeksekusi node \textit{Print String} yang terhubung di alur depannya secara langsung. Hasilnya, sebuah teks akan tercetak singkat di ujung layar saat permainan baru saja dimuat.
\end{itemize}

\begin{figure}[H]
    \centering
    \includegraphics[width=0.6\textwidth]{25.png}
    \caption{Implementasi Blueprint: Event Tick Play}
    \label{fig:25}
\end{figure}
Penjelasan :
\begin{itemize}
    \item \textbf{Cara Kerja:} Pada skenario ini, node yang digunakan adalah \textit{Event Tick}. Oleh karena node ini bekerja secara menerus pada setiap waktu perputaran gambar (\textit{frame}), \textit{Print String} dipanggil secara berulang-ulang tanpa henti. Hal ini menyebabkan layar gim dibanjiri oleh banyak baris teks pesan secara terus-menerus selama objek/level ini berjalan.
\end{itemize}

\begin{figure}[H]
    \centering
    \includegraphics[width=0.6\textwidth]{26.png}
    \caption{Implementasi Blueprint: Logika Toggle Output Teks}
    \label{fig:26}
\end{figure}
Penjelasan :
\begin{itemize}
    \item \textbf{Cara Kerja:} Logika ini merupakan sistem saklar percabangan (\textit{toggle}). Saat permainan dimulai, \textit{Event BeginPlay} memberikan izin agar sistem dapat menerima masukan \textit{keyboard} dari pemain (\textit{Enable Input}). \textit{Event Tick} di atasnya secara konstan mencetak teks berdasarkan nilai variabel Boolean \textbf{Is Active} (lewat percabangan \textit{Branch}). Jika nilainya \textit{True}, sistem mencetak \textbf{"INI BENAR"}, sebaliknya jika \textit{False} mencetak \textbf{"INI TIDAK BENAR"}. 
    \item Interaksi terjadi ketika pemain menekan tombol keyboard \textbf{R}. Begitu ditekan, nilai \textbf{Is Active} yang ada disetel menjadi kebalikannya oleh fungsi node \textbf{NOT}. Akibatnya, arus pencetakan teks di layar akan berganti seketika dari pesan "salah" menjadi "benar" (atau sebaliknya), di setiap kali tombol R diklik.
\end{itemize}

\begin{figure}[H]
    \centering
    \includegraphics[width=0.6\textwidth]{27.png}
    \caption{Implementasi Blueprint: Manipulasi Skala Objek Berulang (Pulsate)}
    \label{fig:27}
\end{figure}
Penjelasan :
\begin{itemize}
    \item \textbf{Cara Kerja:} Rangkaian algoritma visual ini menciptakan efek animasi pernapasan/pemuaian pada objek 3D di udara. Melalui siklus \textit{Event Tick}, sistem senantiasa mengekstrak ukuran skala (\textit{Scale 3D}) dari objek, memecahnya, dan membandingkan sumbu X apakah sudah mencapai nilai batas (\textit{Scale Target}) menggunakan operasi logika $\geq$.
    \item Apabila objek masih kecil (belum mencapai batas/\textit{False}), sistem akan menambahkan skalanya sebesar \textbf{0.001} di setiap \textit{frame} sehingga membesar perlahan. Namun, jika ukurannya sudah menyentuh batas (misal \textbf{1.0}), alur berpindah menjadi \textit{True}. Di sini target batas diubah mengecil (menjadi \textbf{0.1}), dan operasi pengubahan skala dipindah menjadi \textit{pengurangan} ukuran sebesar \textbf{0.001} secara berlanjut. Hasil akhirnya objek tersebut memuai besar dan menyusut secara berulang-ulang tanpa putus.
\end{itemize}

\chapter{Daftar Pustaka}

\begin{enumerate}
    % Rujukan utama yang paling penting untuk Blueprint UE5
    \item \label{epic} Epic Games. \textit{Unreal Engine 5 Documentation}. [Online]. Tersedia: \url{https://docs.unrealengine.com/}.
    
    % Buku yang sangat spesifik membahas Node Blueprint UE5
    \item \label{romero} M. Romero. \textit{Blueprints Visual Scripting for Unreal Engine 5}, 3rd ed. Birmingham, UK: Packt Publishing, 2022.
    
    % Buku fundamental untuk pemahaman konsep game engine (seperti Event Tick, Begin Play)
    \item \label{gregory} J. Gregory. \textit{Game Engine Architecture}, 3rd ed. Boca Raton, FL: CRC Press, 2018.
    
    % Buku tambahan untuk logika Blueprint dasar
    \item \label{dorantes} M. Dorantes. \textit{Unreal Engine for Beginners: Learn to create games with Blueprint}. Independently published, 2020.
    
    % Fundamental grafis/rendering untuk engine (opsional tapi bagus untuk memperkuat laporan)
    \item \label{akenine} T. Akenine-Möller, E. Haines, dan N. Hoffman. \textit{Real-Time Rendering}, 4th ed. Boca Raton, FL: A K Peters/CRC Press, 2018.
\end{enumerate}